This paper presents a minimal common framework for describing biological systems in terms of information density \(\I\), information flow \(\J\), and an induced potential \(\PhiI\). These are source-local formal or model quantities with proposed biological readings, not operationally identified observables. In this framework, DNA is represented abstractly through a parameter set \(\theta\) that constrains system dynamics, rather than as an encoded final form; no concrete genotype-to-parameter map is supplied.

On this basis, development and morphology are given a proposed attractor-based interpretation, and life is described through a candidate sustained non-equilibrium region. The existence, stability, and operational biological identification of such regions remain separate questions. The probability of the emergence of life is correspondingly reformulated as a region-entry target rather than a single-point target; no numerical probability is calculated here.

Intelligence is treated only as a forward conceptual structure: a proposed second attractor supported by the life-attractor description. The paper does not introduce a new biological mechanism. It provides a common structural description using a minimal set of dynamical quantities; the present simulation is a qualitative model-local witness and is not tested here against an external biological dataset.
